\chapter{Banco de Dados}
\label{cap:banco-dados-teoria}

\begin{edital}
Modelagem de dados (conceitual, lógica e física); abordagem relacional e
multidimensional; normalização; integridade referencial; metadados;
modelagem dimensional; SQL; DDL; DML; SGBD; propriedades de banco de dados;
NoSQL; banco de dados em memória; \textit{data lakes} e soluções para big
data; dados estruturados e não estruturados; avaliação de modelos de dados;
técnicas de integração e ingestão (ETL/ELT, transferência de arquivos,
integração via base de dados).
\end{edital}

\section{Modelagem em três níveis}

\begin{conceito}[A mesma pergunta, respondida três vezes, com liberdade decrescente]
Modelar um banco é responder três vezes à mesma pergunta de negócio, cada
vez com \textbf{menos} liberdade de escolha. \textit{O que existe no mundo
do negócio?} --- existem servidores, e cada servidor tem dependentes.
\textit{Como isso vira estrutura de dados?} --- vira duas tabelas, ligadas
por uma chave. \textit{Como isso roda neste produto específico?} --- vira
um \texttt{CREATE TABLE} com tipos concretos, um índice na coluna de busca
e uma escolha de armazenamento.

A separação não é burocracia: os três níveis \textbf{mudam em ritmos
diferentes}. A regra “todo dependente pertence a um servidor” sobrevive à
troca do SGBD Oracle pelo PostgreSQL; o tipo \texttt{NUMBER(5)} específico
do Oracle não sobrevive. Quem mistura os níveis (decide tipo de coluna
enquanto ainda discute com o usuário o que é uma entidade) acaba
redesenhando o negócio toda vez que troca de infraestrutura.

\begin{center}
\begin{tabular}{@{}p{2.4cm}p{4.6cm}p{4.6cm}@{}}
\toprule
\textbf{Nível} & \textbf{O que é} & \textbf{Artefato} \\
\midrule
Conceitual & visão de negócio, independe de SGBD & Modelo
Entidade-Relacionamento (MER/DER) \\
Lógico & estrutura no modelo de dados escolhido (relacional, documento...) &
tabelas, chaves, tipos genéricos \\
Físico & implementação no SGBD específico & DDL, índices, particionamento \\
\bottomrule
\end{tabular}
\end{center}

O critério para saber em que nível uma decisão pertence é \textbf{o que já
foi decidido até ali}: o \textbf{conceitual} decide entidades, atributos e
relacionamentos --- não sabe o que é tabela nem chave estrangeira; é o
nível discutido \textbf{com o usuário de negócio}, sem jargão técnico. O
\textbf{lógico} já escolheu o \textit{paradigma} (relacional, documento,
grafo) e traduz o conceitual para ele --- entidades viram tabelas,
relacionamentos viram chaves estrangeiras ou tabelas associativas --- mas
ainda \textbf{não} escolheu o produto/SGBD. O \textbf{físico} já escolheu o
produto e decide o que só existe ali dentro: tipo exato de coluna, índice,
partição, \textit{tablespace} --- é o único nível em que “desempenho” é uma
resposta legítima à pergunta “por que essa decisão”.

A informação \textbf{só se acrescenta} de um nível para o outro: cada nível
carrega tudo que o anterior decidiu e soma o seu. É por isso que a ordem
não é invertível --- não se parte de um índice físico para descobrir qual
é a entidade de negócio.
\end{conceito}

\subsection{Metadados}

\begin{conceito}[Metadado é dado sobre o dado, não o dado]
\textbf{Metadado} descreve a estrutura, o significado, a origem e as
regras de um dado --- nunca o valor dele. Se \texttt{"1988-10-05"} é o
dado, o metadado é “a coluna chama-se \texttt{data\_admissao}, é do tipo
\texttt{DATE}, não aceita nulo, veio do sistema de folha de pagamento e
significa a data de entrada em exercício do servidor”. Divide-se em três
famílias:

\begin{center}
\begin{tabular}{@{}p{2.6cm}p{9cm}@{}}
\toprule
\textbf{Técnico} & o que o SGBD precisa para funcionar: tabelas, colunas,
tipos, chaves, índices, restrições. Mora no \textbf{dicionário de dados}
(catálogo do sistema), consultável por SQL \\
\textbf{De negócio} & o que a \textbf{pessoa} precisa entender: definição
do termo, regra de cálculo, dono do dado (\textit{data owner}), grau de
sigilo. Conteúdo do \textbf{glossário de negócio} \\
\textbf{Operacional} & o que aconteceu com o dado ao longo do tempo: data
da última carga, volume processado, erros, e a \textbf{linhagem}
(\textit{data lineage}) --- de que fonte ele veio e por quais
transformações passou \\
\bottomrule
\end{tabular}
\end{center}

A \textbf{linhagem} é a que mais rende em questão de cenário: quando um
relatório não fecha, é ela que responde “de onde saiu este número, e o que
foi feito com ele no caminho”. Sem metadado operacional, um pipeline de
ETL vira uma caixa-preta --- ninguém consegue dizer qual regra transformou
o quê.
\end{conceito}

\begin{cuidado}
Metadado \textbf{não é} o conteúdo/valor da coluna (“o CPF do servidor” é
dado, não metadado) --- metadado é a \textbf{descrição} desse dado.
Dicionário de dados (catálogo \textbf{técnico}) e glossário de negócio
(definições em linguagem de \textbf{negócio}) não são a mesma coisa.
\end{cuidado}

\subsection{Avaliação de modelos de dados}

\begin{conceito}[Critérios objetivos, não gosto pessoal]
“O modelo está bom?” tem critérios concretos, não é opinião:

\begin{itemize}
\item \textbf{Completude}: o modelo cobre todos os requisitos levantados
--- nenhuma entidade, atributo ou regra ficou de fora.
\item \textbf{Corretude}: o que está desenhado corresponde à realidade do
negócio --- cardinalidades, opcionalidade e domínios refletem o que
realmente acontece.
\item \textbf{Não redundância / grau de normalização}: cada fato mora num
único lugar (ou a redundância é \textbf{intencional} e justificada, como
num Data Warehouse).
\item \textbf{Integridade}: chaves primárias e estrangeiras declaradas,
restrições explícitas no modelo em vez de confiadas ao código da
aplicação.
\item \textbf{Aderência a padrão}: convenção de nomes, notação e tipos
uniformes --- o que permite outra pessoa ler o modelo sem um tradutor.
\item \textbf{Legibilidade e manutenibilidade}: o modelo se entende sem
quem o desenhou por perto, e uma mudança de regra de negócio não obriga a
redesenhar meia base.
\end{itemize}

O corte que decide entre os dois primeiros critérios: \textbf{completude}
pergunta “falta algo?”; \textbf{corretude} pergunta “o que está lá está
certo?”. Um modelo pode ser \textbf{completo e errado} (tem tudo, mas com
uma cardinalidade invertida) ou \textbf{correto e incompleto} (o que está
desenhado está certo, mas faltou modelar o requisito de auditoria).
\end{conceito}

\section{Modelo relacional e integridade}

\begin{conceito}[A ideia central de Codd (1970) e sua consequência inevitável]
O modelo relacional tem uma ideia central: \textbf{todo dado é uma
tabela, e a única forma de ligar duas tabelas é o valor de um campo
coincidir com o valor de outro} --- não há ponteiro, não há endereço de
memória, não há ordem garantida entre linhas. A consequência inevitável:
se o vínculo é por \textbf{valor}, alguém precisa \textbf{garantir que o
valor referenciado exista}, senão a ligação aponta para o nada. É disso
que trata a integridade referencial.

Caso concreto: se \texttt{DEPENDENTE.servidor\_id} vale 47 e não existe
servidor com id 47, aquela linha é lixo --- nenhuma consulta com
\texttt{JOIN} vai trazê-la de volta corretamente, porque não há do outro
lado quem ela referencia. A restrição de integridade referencial é o que
\textbf{impede essa linha de nascer} no banco.
\end{conceito}

\begin{conceito}[Vocabulário básico]
\begin{itemize}
\item \textbf{Chave primária (PK)}: identifica unicamente cada tupla
(linha); não pode ser nula, precisa ser única.
\item \textbf{Chave estrangeira (FK)}: referencia a PK de outra tabela (ou,
em casos recursivos, da mesma tabela).
\item \textbf{Integridade referencial}: toda FK aponta para uma PK que
existe, ou é nula.
\item \textbf{Integridade de entidade}: a PK nunca pode ser nula.
\item \textbf{Álgebra relacional}: seleção ($\sigma$) corresponde ao
\texttt{WHERE} do SQL e filtra \textbf{linhas}; projeção ($\pi$) filtra
\textbf{colunas}.
\end{itemize}
\end{conceito}

\begin{conceito}[As ações referenciais: o que fazer quando o “pai” é apagado]
A pergunta que \texttt{ON DELETE} responde: apagaram o servidor 47, e
existem dependentes apontando para ele. O que acontece com os
“dependentes órfãos”?

\begin{center}
\begin{tabular}{@{}p{2.8cm}p{9.2cm}@{}}
\toprule
\textbf{Ação} & \textbf{O que o SGBD faz} \\
\midrule
\texttt{CASCADE} & apaga (ou atualiza) os filhos junto --- o dependente
some com o servidor \\
\texttt{SET NULL} & mantém o filho e zera a FK --- exige que a coluna
aceite \texttt{NULL} \\
\texttt{SET DEFAULT} & mantém o filho e aponta a FK para um valor padrão,
que precisa existir \\
\texttt{RESTRICT} & \textbf{barra} a operação de imediato, se existirem
filhos \\
\texttt{NO ACTION} & barra também, mas a checagem só ocorre no \textbf{fim
do comando} --- permite que um passo intermediário viole a restrição,
desde que o estado final esteja íntegro \\
\bottomrule
\end{tabular}
\end{center}

\texttt{RESTRICT} e \texttt{NO ACTION} são quase idênticos: os dois
recusam a operação, e a diferença é só \textbf{quando} a verificação
acontece. \texttt{CASCADE} é o oposto dos dois --- não recusa nada, ele
propaga a mudança. Uma consequência prática: \texttt{CASCADE} numa tabela
central pode apagar em silêncio uma quantidade grande de dados a partir de
um único \texttt{DELETE}.
\end{conceito}

\subsection{Do MER ao relacional: como cada cardinalidade vira tabela}

\begin{conceito}
\begin{center}
\begin{tabular}{@{}p{2.2cm}p{9.8cm}@{}}
\toprule
\textbf{Cardinalidade} & \textbf{Implementação no modelo relacional} \\
\midrule
1:1 & FK em uma das duas tabelas (de preferência no lado de participação
obrigatória), com restrição de unicidade \\
1:N & \textbf{FK no lado N} --- a tabela “muitos” guarda a chave da tabela
“um” \\
\textbf{N:M} & \textbf{tabela associativa} (intermediária), com as FKs das
duas pontas formando a chave primária composta \\
\bottomrule
\end{tabular}
\end{center}
\end{conceito}

\begin{exemplo}[Por que N:M precisa de uma terceira tabela]
Um \texttt{ALUNO} pode se matricular em vários \texttt{CURSO}s, e um
\texttt{CURSO} tem vários \texttt{ALUNO}s. Se a FK ficasse dentro de
\texttt{ALUNO} (apontando para um curso só), cada aluno só poderia estar
em um curso --- perderia a possibilidade de matricular-se em vários. A
solução é criar \texttt{MATRICULA(aluno\_id, curso\_id)}, com as duas FKs
formando a chave primária composta --- cada linha dessa tabela representa
uma combinação aluno-curso específica.
\end{exemplo}

\begin{cuidado}
N:M \textbf{nunca} se resolve com FK direta numa das tabelas (isso só
atende 1:N) nem com coluna multivalorada dos dois lados --- coluna
multivalorada viola a 1FN (Seção~\ref{sec:normalizacao}).
\end{cuidado}

\subsection{Entidade fraca x entidade associativa}

\begin{conceito}
\textbf{Entidade fraca}: não se identifica sozinha. Sua chave é a chave da
\textbf{entidade proprietária} (forte) somada a uma \textbf{chave parcial}
(discriminador) --- por exemplo, \texttt{DEPENDENTE} só é identificado
dentro do cadastro de um \texttt{SERVIDOR} específico, pelo nome. O
relacionamento com o proprietário é \textbf{identificador}, e a existência
da entidade fraca depende da entidade forte.

\textbf{Entidade associativa}: conceito diferente --- é o
\textbf{relacionamento N:M promovido a entidade}, usado quando esse
relacionamento precisa, ele próprio, se relacionar com uma terceira
entidade ou ganhar atributos próprios (por exemplo, \texttt{MATRICULA} da
seção anterior, se ganhar o atributo \texttt{nota}, vira uma entidade
associativa).
\end{conceito}

\begin{cuidado}
Entidade fraca \textbf{não dispensa} a chave do proprietário --- é
exatamente o oposto: sem ela, a entidade fraca não se identifica sozinha.
Ter atributos próprios \textbf{não} torna uma entidade “forte”; o que
define força é conseguir se identificar sozinha, sem depender da chave de
outra entidade.
\end{cuidado}

\section{Normalização}
\label{sec:normalizacao}

\begin{conceito}[O problema real que a normalização resolve: anomalias]
Normalizar não é obediência a uma regra estética --- resolve um problema
concreto chamado \textbf{anomalia}, que aparece sempre que uma tabela
guarda \textbf{dois assuntos ao mesmo tempo}.
\end{conceito}

\begin{conceito}[Dependência funcional --- a ferramenta que sustenta tudo]
Escreve-se \textbf{A $\to$ B} e lê-se “\textbf{A determina B}”: dado um
valor de A, existe \textbf{um único} valor de B possível. \texttt{CPF
$\to$ nome} é dependência funcional --- um CPF corresponde a um único
nome. \texttt{nome $\to$ CPF} \textbf{não} é --- homônimos existem, o mesmo
nome pode ter CPFs diferentes.

Toda forma normal é uma restrição sobre \textbf{quais} dependências
funcionais uma tabela tem o direito de conter. A chave determina todo o
resto --- isso é literalmente o que significa “ser chave”. O que a
normalização proíbe, forma a forma, é que \textbf{algo além da chave
inteira} determine outra coisa:

\begin{itemize}
\item determinar a partir de \textbf{parte} da chave (não da chave
inteira) $\to$ viola a \textbf{2FN} (dependência parcial) --- só é
possível quando a chave é composta, por isso a 2FN não tem o que fazer
numa tabela de chave simples;
\item determinar a partir de um atributo \textbf{não-chave} (um atributo
comum determinando outro) $\to$ viola a \textbf{3FN} (dependência
transitiva);
\item determinar sendo algo que \textbf{não é chave candidata} $\to$
viola a \textbf{BCNF} (forma normal de Boyce-Codd).
\end{itemize}

As três são a mesma frase ficando progressivamente mais rigorosa. Por
isso a ordem é cumulativa: quem está em 3FN já está automaticamente em
2FN e em 1FN.
\end{conceito}

\begin{center}
\begin{tabular}{@{}p{2cm}p{9.4cm}@{}}
\toprule
\textbf{FN} & \textbf{Elimina} \\
\midrule
1FN & atributos multivalorados / não atômicos (cada célula, um valor só) \\
2FN & dependência parcial da chave (só é relevante com PK composta) \\
3FN & dependência transitiva (atributo não-chave que depende de outro
atributo não-chave) \\
BCNF & anomalia residual quando existe mais de uma chave candidata
sobreposta \\
\bottomrule
\end{tabular}
\end{center}

Normalizar \textbf{reduz redundância}; \textbf{desnormalizar de
propósito} (caso típico: Data Warehouse, Capítulo \ref{cap:bi-teoria}) é
aceitável quando o objetivo é desempenho de \textbf{leitura} --- nunca de
escrita.

\subsection{Decomposição passo a passo: um exemplo completo}

\begin{exemplo}[De uma tabela com anomalias a três tabelas em 3FN]
Uma tabela de matrículas mal projetada, guardando aluno, curso, nota,
nome e departamento do aluno, e chefe do departamento, tudo junto:

\begin{center}
\small
\begin{tabular}{@{}llllll@{}}
\toprule
\textbf{aluno} & \textbf{curso} & \textbf{nota} & \textbf{nome\_aluno} &
\textbf{depto} & \textbf{chefe\_depto} \\
\midrule
101 & BD01 & 8,5 & Ana   & Computação & Prof. Silva \\
101 & ES02 & 7,0 & Ana   & Computação & Prof. Silva \\
102 & BD01 & 9,0 & Bruno & Computação & Prof. Silva \\
\bottomrule
\end{tabular}
\end{center}

A chave é composta: \textbf{(aluno, curso)} --- é o par que identifica uma
matrícula específica. As dependências funcionais presentes:

\begin{center}
\begin{tabular}{@{}p{5cm}p{7.4cm}@{}}
\toprule
\textbf{Dependência} & \textbf{Status} \\
\midrule
(aluno, curso) $\to$ nota & legítima: depende da chave \textbf{inteira} \\
aluno $\to$ nome\_aluno, depto & \textbf{parcial}: depende de só
\textbf{metade} da chave (viola 2FN) \\
depto $\to$ chefe\_depto & \textbf{transitiva}: um atributo não-chave
determina outro (viola 3FN) \\
\bottomrule
\end{tabular}
\end{center}

\textbf{As três anomalias que isso causa}: (1) \textbf{de atualização}
--- o chefe da Computação muda; ele está repetido em toda linha de todo
aluno do departamento, e se um \texttt{UPDATE} esquecer uma linha, a
tabela passa a afirmar duas coisas contraditórias ao mesmo tempo; (2)
\textbf{de inserção} --- criou-se um departamento novo, ainda sem aluno
nenhum; não há onde guardá-lo, porque a chave exige (aluno, curso) e não
existe aluno; (3) \textbf{de exclusão} --- o aluno 102 trancou e sua
única matrícula foi apagada; junto com ela sumiu o próprio registro de que
Bruno existe.

\textbf{Decompondo, uma forma normal por vez.} Já está em \textbf{1FN}
(toda célula tem um valor só). Para \textbf{2FN}, elimina-se a
dependência parcial: nome e departamento do aluno não dependem de
\textit{curso}, dependem só do \textit{aluno} --- separa-se uma tabela:

\begin{center}
\texttt{MATRICULA(aluno, curso, nota)}\\
\texttt{ALUNO(aluno, nome\_aluno, depto, chefe\_depto)}
\end{center}

Para \textbf{3FN}, elimina-se a transitiva: dentro de \texttt{ALUNO}, o
chefe não depende diretamente do aluno, depende do \textit{departamento},
que por sua vez depende do aluno --- separa-se mais uma:

\begin{center}
\texttt{MATRICULA(aluno, curso, nota)}\\
\texttt{ALUNO(aluno, nome\_aluno, depto)}\\
\texttt{DEPARTAMENTO(depto, chefe\_depto)}
\end{center}

Cada fato passou a morar num único lugar, e as três anomalias
desaparecem junto: trocar o chefe vira um \texttt{UPDATE} numa única
linha; cadastrar um departamento novo não exige inventar um aluno; apagar
a última matrícula de alguém não apaga o registro do aluno.
\end{exemplo}

\begin{regrapratica}[Achar a chave a partir das dependências funcionais]
A FGV costuma cobrar assim: “considere $R(A,B,C,D,E)$ e $F = \{A \to B,
BC \to DE, \ldots\}$; qual a forma normal mais alta?”. O procedimento é
mecânico: parta dos atributos que \textbf{nunca aparecem do lado direito}
de nenhuma dependência --- eles obrigatoriamente fazem parte de toda chave
candidata. Feche esse conjunto aplicando as dependências até parar de
crescer; se ele alcançar todos os atributos da relação, é chave. Só então
classifique cada dependência restante como parcial, transitiva, ou fora de
chave candidata --- a forma normal mais alta é a última que \textbf{nenhuma}
dependência viola.
\end{regrapratica}

\section{SQL: DDL x DML x DCL x TCL}

\begin{conceito}[A divisão é sobre o que cada grupo faz com a transação]
DDL muda a \textbf{estrutura} do banco e, na maioria dos SGBD, confirma
sozinha (\textit{auto-commit}) --- não se desfaz um \texttt{DROP TABLE}
com \texttt{ROLLBACK}. DML muda os \textbf{dados} dentro de uma
transação, e por isso é desfazível. É essa diferença de comportamento
transacional --- não a semelhança do efeito visível --- que separa
\texttt{TRUNCATE} (DDL) de \texttt{DELETE} (DML), mesmo que os dois
“esvaziem a tabela”.

\begin{center}
\begin{tabular}{@{}p{2.6cm}p{4cm}p{4.5cm}@{}}
\toprule
\textbf{Sublinguagem} & \textbf{Comandos} & \textbf{Observação} \\
\midrule
DDL (definição) & CREATE, ALTER, \textbf{DROP, TRUNCATE} & TRUNCATE é DDL
(não DML), \textit{auto-commit} \\
DML (manipulação) & SELECT, INSERT, UPDATE, DELETE & DELETE é DML, aceita
WHERE e é desfazível com ROLLBACK \\
DCL (controle) & GRANT, REVOKE & gerencia permissões \\
TCL (transação) & COMMIT, ROLLBACK, SAVEPOINT & controla o ciclo de vida
da transação \\
\bottomrule
\end{tabular}
\end{center}
\end{conceito}

\begin{conceito}[A ordem lógica de execução, e o que ela explica]
Você escreve o \texttt{SELECT} primeiro na sintaxe, mas o banco o resolve
quase por último:

\begin{center}
\texttt{FROM} $\to$ \texttt{WHERE} $\to$ \texttt{GROUP BY} $\to$
\texttt{HAVING} $\to$ \texttt{SELECT} $\to$ \texttt{ORDER BY}
\end{center}

Duas consequências práticas que decorrem diretamente dessa ordem:

\begin{itemize}
\item \textbf{O \texttt{WHERE} não enxerga um apelido definido no
\texttt{SELECT}}. Em \texttt{SELECT preco*2 AS dobro ... WHERE dobro >
100}, o \texttt{WHERE} roda \textbf{antes} de \texttt{dobro} existir --- o
banco acusa erro de coluna inexistente. Já \texttt{ORDER BY dobro}
funciona normalmente, porque \texttt{ORDER BY} roda depois do
\texttt{SELECT}.
\item \textbf{\texttt{WHERE} filtra linha; \texttt{HAVING} filtra
grupo}. O \texttt{WHERE} acontece \textbf{antes} do \texttt{GROUP BY},
quando nenhum grupo existe ainda; o \texttt{HAVING} acontece depois, e por
isso é o único que pode comparar com o resultado de \texttt{COUNT(*)} ou
\texttt{SUM()}.
\end{itemize}
\end{conceito}

\begin{cuidado}
\texttt{UNION} remove duplicatas; \texttt{UNION ALL} mantém todas as
linhas, inclusive repetidas. \texttt{JOIN} sem cláusula \texttt{ON}
produz um \textbf{produto cartesiano} (cada linha de uma tabela combinada
com todas as linhas da outra).
\end{cuidado}

\subsection{VIEW, GRANT e controle de transação}

\begin{conceito}
\textbf{VIEW (visão)}: uma \textbf{tabela virtual} --- é a consulta
guardada com um nome, não os dados copiados fisicamente. Consultar a
visão executa o \texttt{SELECT} que está por trás dela, então o resultado
é \textbf{sempre atual}. Serve para simplificar uma consulta complexa e
para \textbf{restringir acesso} (a visão expõe só as colunas/linhas
permitidas, e o \texttt{GRANT} é dado sobre a visão, não sobre a tabela
inteira). Uma visão \textbf{simples} costuma aceitar atualização; uma
visão com \texttt{JOIN}, agregação ou \texttt{DISTINCT}, geralmente não. A
\textbf{\textit{materialized view}} é a exceção: ela \textbf{sim}
armazena o resultado fisicamente em disco, e precisa ser
\textbf{atualizada} (\textit{refresh}) periodicamente para refletir
mudanças na origem.

\textbf{GRANT/REVOKE (DCL)}: \texttt{GRANT SELECT, INSERT ON tabela TO
usuario} concede privilégios; \texttt{REVOKE} os retira. Com
\texttt{WITH GRANT OPTION}, quem recebeu o privilégio pode
\textbf{repassá-lo} a terceiros.

\textbf{SAVEPOINT e ROLLBACK TO (TCL)}: \texttt{SAVEPOINT nome} marca um
ponto \textbf{dentro} de uma transação em andamento; \texttt{ROLLBACK TO
nome} desfaz só o que veio \textbf{depois} daquela marca, e a transação
\textbf{continua aberta} --- não equivale a um \texttt{COMMIT} nem a um
\texttt{ROLLBACK} total. É um desfazer parcial.
\end{conceito}

\begin{cuidado}
Uma visão comum \textbf{não armazena} dados fisicamente (isso é exclusivo
da \textit{materialized view}). \texttt{GRANT} é \textbf{DCL}, não DDL
nem DML. \texttt{ROLLBACK TO} \textbf{não encerra} a transação nem desfaz
tudo --- ele volta ao \textit{savepoint} indicado, e a transação segue
viva.
\end{cuidado}

\subsection{Notação Crow's Foot e chave surrogada}

\begin{conceito}[Crow's Foot: dois símbolos por ponta, dois significados]
Cada ponta de uma linha de relacionamento traz \textbf{dois} símbolos, e a
posição decide o papel: o que \textbf{encosta na entidade} é o
\textbf{máximo} (a cardinalidade propriamente dita --- um ou muitos); o
símbolo \textbf{mais afastado} é o \textbf{mínimo} (a modalidade --- zero
ou um).

\begin{center}
\begin{tabular}{@{}p{4.6cm}p{7.4cm}@{}}
\toprule
\textbf{Símbolo} & \textbf{Leitura} \\
\midrule
Traço (barra) & exatamente \textbf{um} (mínimo 1 / máximo 1) \\
Círculo & \textbf{zero} --- opcional \\
Pé de galinha (três riscos) & \textbf{muitos} \\
Círculo + pé de galinha & \textbf{zero ou muitos} \\
Traço + pé de galinha & \textbf{um ou muitos} (obrigatório) \\
\bottomrule
\end{tabular}
\end{center}

Tradução direta para DDL: o lado com \textbf{pé de galinha} é o lado
\textbf{N}, e é nele que entra a \textbf{FK}; o círculo indica
participação \textbf{opcional} (a FK aceita \texttt{NULL}); o traço
indica participação \textbf{obrigatória} (\texttt{NOT NULL}).
\end{conceito}

\begin{conceito}[Chave surrogada x chave natural]
\textbf{Chave surrogada} (substituta/artificial): chave primária
\textbf{sem significado de negócio}, gerada pelo sistema (sequência,
\textit{identity}, UUID) --- em oposição à \textbf{chave natural}, que
vem do próprio domínio (CPF, matrícula). Vantagens da surrogada: é
\textbf{estável} (um dado de negócio pode mudar --- CPF corrigido,
matrícula reemitida --- sem exigir alterar a chave e todas as FKs que a
referenciam), curta e uniforme para \textit{join}. É padrão em Data
Warehouse, onde também permite guardar \textbf{versões históricas} da
mesma entidade de negócio (dimensão de mudança lenta, Capítulo
\ref{cap:bi-teoria}).
\end{conceito}

\begin{cuidado}
Chave surrogada \textbf{não substitui} a necessidade da chave natural na
tabela --- o campo de negócio continua presente, em geral com restrição
\texttt{UNIQUE}. Em Crow's Foot, a inversão comum é ler a cardinalidade
do lado errado da linha --- o símbolo vale para a entidade que ele
\textbf{toca}, não para a entidade do outro lado.
\end{cuidado}

\subsection{Gatilho x procedimento armazenado}

\begin{conceito}
\begin{center}
\begin{tabular}{@{}p{2.4cm}p{4.8cm}p{4.8cm}@{}}
\toprule
& \textbf{Gatilho (\textit{trigger})} & \textbf{Procedimento armazenado} \\
\midrule
Quem dispara & o \textbf{próprio SGBD}, em resposta a um \textbf{evento}
de dados & a \textbf{aplicação} (ou outro código), por chamada explícita \\
Como se invoca & não se invoca --- é automático & \texttt{CALL}/
\texttt{EXECUTE} \\
Amarrado a & uma tabela + evento (INSERT, UPDATE, DELETE) & nada
específico; é código nomeado e reutilizável \\
Parâmetros & não recebe & recebe (e pode retornar) \\
Bom para & auditoria, log, regra que \textbf{não pode depender} de a
aplicação lembrar de chamá-la & rotina de negócio chamada sob demanda \\
\bottomrule
\end{tabular}
\end{center}
Palavra-chave para reconhecer o gatilho num enunciado: “\textbf{sem
depender de a aplicação lembrar}”.
\end{conceito}

\section{Propriedades ACID}

\begin{conceito}[As quatro letras, na transferência de R\$ 100 de A para B]
Uma transação é um conjunto de operações que o banco trata como
\textbf{uma só}. Debitar R\$ 100 da conta A e creditar R\$ 100 na conta B
são dois \texttt{UPDATE}s, e nenhum estado intermediário entre eles é
aceitável.

\begin{itemize}
\item \textbf{A}tomicidade --- \textit{tudo ou nada}. Se o servidor cair
entre o débito e o crédito, o débito é desfeito. Nunca existe “saiu de A
e não chegou em B”. Quem entrega: \texttt{ROLLBACK} e o log de
transações.
\item \textbf{C}onsistência --- \textit{estado válido antes, estado
válido depois}. A soma dos saldos não mudou, e nenhuma restrição
declarada (\texttt{CHECK}, chave estrangeira, \texttt{UNIQUE}) foi
violada. Consistência é sobre \textbf{as regras do banco}, não sobre
concorrência entre transações.
\item \textbf{I}solamento --- \textit{uma transação não enxerga o meio da
outra}. Se um extrato rodar no exato instante da transferência, ele não
pode ver o dinheiro debitado e ainda não creditado. É a letra mais
cobrada porque é a única que tem \textbf{graus} (Seção seguinte).
\item \textbf{D}urabilidade --- \textit{confirmado é para sempre}. Depois
do \texttt{COMMIT}, o efeito sobrevive a uma queda de energia. Quem
entrega: a gravação do log em disco \textbf{antes} da confirmação (WAL,
\textit{write-ahead log}).
\end{itemize}
\end{conceito}

\begin{cuidado}
Consistência x isolamento se confundem: a primeira olha para as
\textbf{regras declaradas} do banco; a segunda olha para \textbf{outras
transações} concorrentes. Não confundir com o “C” do teorema CAP
(Seção~\ref{sec:cap-teoria}) --- lá, consistência significa que todas as
réplicas mostram o mesmo dado, um conceito diferente.
\end{cuidado}

\section{Concorrência: níveis de isolamento}

\begin{conceito}[Quatro níveis, cada um admitindo certos fenômenos]
O “I” do ACID não é tudo-ou-nada: o padrão SQL define quatro níveis, e
cada um \textbf{admite} certos fenômenos de concorrência. Quanto maior o
isolamento, \textbf{menor} a concorrência possível --- é uma troca
(\textit{trade-off}), não um ganho de graça.

\textbf{Os três fenômenos:} \textbf{leitura suja} (\textit{dirty read}):
lê um dado alterado por outra transação que \textbf{ainda não confirmou}
(e pode desfazer depois). \textbf{Leitura não repetível}
(\textit{non-repeatable read}): relê a \textbf{mesma linha} e o valor
mudou, porque outra transação a atualizou e confirmou no meio do caminho.
\textbf{Leitura fantasma} (\textit{phantom read}): relê a \textbf{mesma
faixa} de resultados e aparecem \textbf{linhas novas}, inseridas e
confirmadas por outra transação.

\begin{center}
\begin{tabular}{@{}p{3.6cm}p{2.6cm}p{2.6cm}p{2.6cm}@{}}
\toprule
\textbf{Nível} & \textbf{Leitura suja} & \textbf{Não repetível} &
\textbf{Fantasma} \\
\midrule
READ UNCOMMITTED & admite & admite & admite \\
READ COMMITTED & \textbf{impede} & admite & admite \\
REPEATABLE READ & impede & \textbf{impede} & admite \\
SERIALIZABLE & impede & impede & \textbf{impede} \\
\bottomrule
\end{tabular}
\end{center}
\end{conceito}

\begin{cuidado}
READ UNCOMMITTED é o \textbf{mais permissivo} (não o mais restritivo) ---
é justamente o que deixa ler dado ainda não confirmado. REPEATABLE READ
elimina a leitura \textbf{não repetível}, mas \textbf{não} elimina a
fantasma (que só é eliminada no SERIALIZABLE). Quanto mais alto o
isolamento, \textbf{menor} a concorrência --- nunca o contrário. Para
diferenciar os dois últimos fenômenos: linha que \textbf{muda de valor}
$\to$ não repetível; linha \textbf{nova que aparece} $\to$ fantasma.
\end{cuidado}

\section{Relacional x multidimensional: OLTP x OLAP}

\begin{conceito}[Duas cargas de trabalho opostas exigem dois modelos opostos]
O sistema que registra uma venda precisa gravar \textbf{uma linha por
vez, muitas vezes por segundo}, sem perder nada --- a normalização serve
bem a ele, porque cada fato mora num único lugar e a escrita fica barata.
O sistema que responde “quanto vendemos por região no último trimestre”
precisa \textbf{varrer milhões de linhas e somar} --- e aí a
normalização vira inimiga, porque cada \texttt{JOIN} a mais é trabalho a
mais numa consulta que já é pesada.

Rodar as duas cargas na mesma base é o erro clássico: o relatório trava a
produção. Separar em OLTP e OLAP é a resposta --- e é isso que justifica
o Data Warehouse existir como uma cópia desnormalizada dos dados, em vez
de ser um capricho de arquitetura.

\begin{center}
\begin{tabular}{@{}p{2.6cm}p{4.4cm}p{4.4cm}@{}}
\toprule
& \textbf{OLTP (transacional)} & \textbf{OLAP (analítico)} \\
\midrule
Uso & operações do dia a dia & apoio à decisão / análise \\
Escrita/leitura & muitas escritas curtas & leitura pesada, agregação \\
Modelo & relacional normalizado & multidimensional (cubos) \\
Estrutura & tabelas & dimensões e métricas/fatos \\
\bottomrule
\end{tabular}
\end{center}
\end{conceito}

\section{NoSQL e novos armazenamentos}

\begin{conceito}[Quatro famílias de NoSQL, e o que cada uma otimiza]
NoSQL prioriza alta disponibilidade e \textbf{escalabilidade horizontal},
e costuma relaxar as garantias ACID em favor do modelo \textbf{BASE}
(\textit{Basically Available, Soft state, Eventual consistency}).

\begin{itemize}
\item \textbf{Chave-valor} (Redis): busca extremamente rápida por uma
chave, sem estrutura interna imposta ao valor.
\item \textbf{Documento} (MongoDB): documentos semiestruturados (tipo
JSON), bom para dados que variam de forma entre registros.
\item \textbf{Colunar} (Cassandra): otimizado para escrita massiva e
leitura por coluna em grandes volumes.
\item \textbf{Grafo} (Neo4j): nós e arestas, otimizado para percorrer
\textbf{relacionamentos} complexos (rede social, detecção de fraude).
\end{itemize}

Banco NoSQL \textbf{não} é “sempre grafo” e não segue ACID estritamente
por padrão --- e não substitui um ERP/CRM fortemente transacional sem
uma análise cuidadosa do que se perde.
\end{conceito}

\subsection{Teorema CAP}
\label{sec:cap-teoria}

\begin{conceito}[Por que o CAP obriga a escolher: o cenário do cabo cortado]
Duas réplicas do mesmo banco, uma no Rio e outra em Natal. O cabo de rede
entre elas cai --- isso é a \textbf{partição} (P). Chega uma escrita no
Rio. O sistema tem exatamente duas saídas:

\begin{itemize}
\item \textbf{Aceitar a escrita}: o Rio responde ao usuário, mas Natal
segue com o dado antigo --- as réplicas divergiram. Abriu-se mão de
\textbf{C}onsistência para manter \textbf{A}vailability: perfil
\textbf{AP}.
\item \textbf{Recusar a escrita} (ou travar até o cabo voltar): ninguém
lê dado divergente, mas o serviço ficou indisponível para aquela
operação. Abriu-se mão de \textbf{A} para manter \textbf{C}: perfil
\textbf{CP}.
\end{itemize}

Não existe terceira saída --- é por isso que é um teorema, não uma
recomendação. “Escolha 2 de 3” é uma simplificação didática: numa rede
real, a partição \textbf{não é opcional} (cabo cai eventualmente); P está
sempre presente, e a escolha verdadeira é sempre entre \textbf{C ou A}. Um
banco monolítico rodando num servidor só \textbf{não é “CA”} --- ele
simplesmente não é um sistema distribuído, e o teorema não se aplica a
ele.
\end{conceito}

\begin{regrapratica}
Pergunte \textbf{o que dói mais} no cenário descrito. Saldo bancário,
sentença judicial, registro oficial, estoque de item único $\to$ dado
errado é inaceitável $\to$ \textbf{CP}. Carrinho de compras, curtida,
\textit{feed} social, telemetria $\to$ ficar fora do ar é o inaceitável, e
a divergência se resolve depois $\to$ \textbf{AP}.
\end{regrapratica}

\subsection{Data Lake x Data Warehouse x Lakehouse}

\begin{conceito}
\textbf{Data Lake}: guarda dado \textbf{bruto} (\textit{schema-on-read}
--- a estrutura só é imposta na hora de ler, não na hora de gravar),
aceita dado estruturado e não estruturado. \textbf{Data Warehouse}: guarda
dado \textbf{tratado e modelado} (\textit{schema-on-write} --- a estrutura
é definida antes de gravar). \textbf{Lakehouse}: combina os dois --- a
flexibilidade de ingestão do lake com as garantias de estrutura e
qualidade do warehouse.
\end{conceito}

\section{ETL x ELT}

\begin{conceito}[A posição do “T” muda tudo]
\begin{center}
\begin{tabular}{@{}p{2.6cm}p{4.3cm}p{4.3cm}@{}}
\toprule
& \textbf{ETL} & \textbf{ELT} \\
\midrule
Ordem & Extrai $\to$ \textbf{Transforma} $\to$ Carrega & Extrai $\to$
Carrega $\to$ \textbf{Transforma} \\
Transformação & antes de carregar, fora do destino & depois de carregar,
dentro do próprio destino \\
Bom para & Data Warehouse tradicional & destino com muito poder de
processamento (nuvem, big data) \\
\bottomrule
\end{tabular}
\end{center}
ELT aproveita o processamento do próprio destino, o que o torna melhor
para \textbf{grande} volume (não o contrário); ETL faz mais sentido
quando a transformação precisa ocorrer antes/fora do destino, por
limitação de processamento lá.
\end{conceito}

\section{Índices}

\begin{conceito}[O que um índice é, e por que existem tipos diferentes]
Um índice é uma \textbf{estrutura separada da tabela}, que guarda os
valores de uma coluna já \textbf{organizados}, junto com o endereço da
linha correspondente. Não muda a tabela em si: acrescenta um caminho de
acesso a mais. A troca é sempre a mesma --- \textbf{leitura mais rápida,
escrita mais lenta} (todo \texttt{INSERT} precisa atualizar o índice
também) e mais espaço em disco. Índice não é grátis, e “criar índice em
tudo” nunca é a resposta certa.

\begin{itemize}
\item \textbf{B+ Tree}: árvore balanceada com os valores \textbf{ordenados}
nas folhas, encadeadas entre si. Como está ordenado, atende faixa
(\texttt{BETWEEN}, \texttt{>}, \texttt{<}), \texttt{ORDER BY} \textbf{e}
igualdade. É o padrão em todo SGBD, ideal para colunas de \textbf{alta
cardinalidade} (muitos valores distintos: CPF, data, valor monetário).
\item \textbf{Hash}: aplica uma função ao valor e vai direto ao “balde”
correspondente. Ótimo para \texttt{=} e \textbf{inútil} para faixa --- o
hash embaralha a ordem de propósito, então “entre 10 e 20” não significa
nada dentro dele.
\item \textbf{Bitmap}: um vetor de bits por valor distinto. Só compensa em
\textbf{baixa cardinalidade} (ex.: coluna ativo/inativo, ou UF com 27
valores), porque o número de vetores é o número de valores distintos.
Combina vários filtros com operações lógicas de bits rapidamente ---
excelente em consulta analítica, ruim em tabela com muita escrita
concorrente.
\end{itemize}

As duas perguntas que decidem o tipo certo: \textit{a consulta pede faixa
ou igualdade?} e \textit{quantos valores distintos a coluna tem?}
\end{conceito}

\section{Big Data: os “V”s}

\begin{conceito}[O trio clássico, e o que veio depois]
O trio \textbf{original} (Gartner, 2001) é \textbf{Volume, Velocidade e
Variedade}. \textbf{Veracidade} e \textbf{Valor} são \textbf{extensões}
posteriores (por isso os modelos “4V” e “5V”), assim como
\textbf{Variabilidade}. Saber qual é o trio original importa porque a
pergunta clássica de prova troca um membro do trio por um V de extensão
(“volume, velocidade e valor” ou “volume, variedade e veracidade” soam
plausíveis mas não são o trio original).
\end{conceito}
